\documentclass[10pt,a4paper]{amsart}
\usepackage[margin=19mm]{geometry}
\usepackage{amsmath,amssymb,mathtools}
\usepackage[scr=rsfs,cal=euler]{mathalfa}
\usepackage{xfrac}
\usepackage[unicode,hidelinks,bookmarks=false]{hyperref}
\newcommand{\ie}{i.e.\ }
\newcommand{\cf}{cf.\ }
\newcommand{\resp}{respectively\ }
\newcommand{\Subsection}{\subsection}
\providecommand{\nobreakdash}{\nobreak\hbox{-}}
\setlength{\emergencystretch}{2em}
\multlinegap0pt
\usepackage{bm}
\usepackage{bbm}
\usepackage{dsfont}
\usepackage{xspace}
\usepackage{cancel}
\usepackage{mathtools}
\usepackage{amsthm}
\makeatletter
\@ifundefined{lemma}{\newtheorem{lemma}{Lemma}}{}
\makeatother
\newcommand\norm[1]{\left\lVert#1\right\rVert}
\newcommand{\Div}{\operatorname{div}}
\newcommand{\Grad}{\nabla}
\newcommand{\N}{\mathbb{N}}
\newcommand{\PVh}{\mathcal{P}_{\Vh}}
\newcommand{\R}{\mathbb{R}}
\newcommand{\Vh}{\mathbb{V}_h}
\newcommand{\dom}{\Omega}
\newcommand{\hu}{\widehat{u}}
\newcommand{\tu}{\widetilde{u}}
\newcommand{\weak}{\rightharpoonup}
\title[A convergence proof for a finite element discretization of C]{A convergence proof for a finite element discretization of Chorin's projection method of the incompressible Navier-Stokes equations}
\author{Franziska Weber}
\date{Source: arXiv:2508.13416v2 (CC BY 4.0)}
\begin{document}
\maketitle

\noindent\emph{Source and licence.} A convergence proof for a finite element discretization of Chorin's projection method of the incompressible Navier-Stokes equations. Version arXiv:2508.13416v2 (CC BY 4.0), distributed under the Creative Commons Attribution 4.0 International licence, as recorded in the arXiv licence metadata for this exact version and shown on the abstract page https://arxiv.org/abs/2508.13416v2. Full rights notice: licenses/arxiv-2508.13416-CC-BY-4.0.txt.\par\smallskip

\noindent\textit{Editorial scope.} Counted unit: the Lemma whose source label is \texttt{lem:aubinlionsmodified} in the article (source0.tex lines 826--841, source label \texttt{lem:aubinlionsmodified}), delivered together with the article's own complete original proof of it (source0.tex lines 842--911, one \texttt{proof} environment of 70 lines). No mathematics is written, repaired or re-derived: statement and proof are the article's own text line for line. Required context retained uncounted: lem:samelimits (lines 645--657); lem:divfree2 (lines 689--695); lem:lionslike (lines 752--758). Every definition, display and label of those lines is retained unchanged. Omitted from this package and not redistributed: 1--644, 658--688, 696--751, 759--825, 912--1326. Nothing inside a retained excerpt is omitted or altered: every retained body is delivered exactly as the article's lines are written, so no omission receipt of any kind accompanies this package. Citations: the retained excerpts carry no citation command at all, so no bibliography entry of the article is transcribed and no reference is redistributed. Reference adaptation: none was needed and none was made; every reference inside the delivered unit resolves inside it, so no unresolved reference or citation is produced. Printed slips kept literal: no printed word, symbol, hypothesis or number of the counted statement or of its proof was altered. Layout and class: the article's own class is \texttt{amsart} with packages including color, geometry, subfig, multicol, listings, amsmath, amsthm, amssymb, wasysym, verbatim, bbm, xcolor, graphics, url; the wrapper is the lane's pinned \texttt{amsart} editorial wrapper at 10pt on A4, which restates the theorem-like environments the retained text needs and adds the article's own macro definitions that the retained text needs (\texttt{\textbackslash Div}, \texttt{\textbackslash Grad}, \texttt{\textbackslash N}, \texttt{\textbackslash PVh}, \texttt{\textbackslash R}, \texttt{\textbackslash Vh}, \texttt{\textbackslash dom}, \texttt{\textbackslash hu}, \texttt{\textbackslash norm}, \texttt{\textbackslash tu}, \texttt{\textbackslash weak}), with extra wrapper packages (bm, bbm, dsfont, xspace, cancel, mathtools). The delivered numbering is the unit's own. Assets: no retained span contains a figure environment, a graphics-inclusion command, a PDF-page inclusion, a table or a TikZ picture; no image file is copied into this package, the package declares no asset dependency and no image file is redistributed with it. Licence: version arXiv:2508.13416v2 of this article is distributed under the Creative Commons Attribution 4.0 International licence, as recorded in the arXiv licence metadata for this exact version and shown on the abstract page https://arxiv.org/abs/2508.13416v2. A full rights notice is delivered at licenses/arxiv-2508.13416-CC-BY-4.0.txt. Characters: every retained excerpt is pure ASCII, so the delivered engine, XeLaTeX with its default Latin Modern text font, reports no missing character.
\begin{lemma}
	\label{lem:samelimits}
	Assume that  $\Delta t = o_{h\to 0}(1)$. Then, we have 
	$u=\bar{u}=\tu=\hu$ a.e. in $[0,T]\times\dom$ and in $L^2([0,T]\times\dom)$, and
	\begin{align}
		\label{eq:uubarconv}
		\lim_{h\to 0}\norm{u_h-\hu_h}_{L^2([0,T]\times\dom)} &= 0,\\
		\label{eq:utildeuhatconv}
		\lim_{h\to 0}\norm{\hu_h-\tu_h}_{L^2([0,T]\times\dom)}& = 0,\\
		\label{eq:utildeubarconv}
		\lim_{h\to 0}\norm{\bar{u}_h-\tu_h}_{L^2([0,T]\times\dom)}& = 0.
	\end{align}	
\end{lemma}

\begin{lemma}\label{lem:divfree2}
	Any weak limit $u$ of the sequence $\{u_h\}_{h>0}$ satisfies for a.e. $t\in [0,T]$,
	\begin{equation*}
		(u,\Grad q) = 0,\quad \forall \, q\in H^1(\dom).
	\end{equation*}
	
\end{lemma}

\begin{lemma}\label{lem:lionslike}
	Denote $Z = (H^k(\dom)\cap H^1_{\Div}(\dom))^*$ for some $k\in \N$. Then for any $\eta>0$ there exists a constant $C_\eta>0$ and $h_\eta>0 $ such that for all $h\leq h_{\eta}$ we have
	\begin{equation}\label{eq:lionslike}
		\norm{\PVh w}_{L^2(\dom)}\leq \eta \norm{w}_{H^1_0(\dom)} + C_\eta \norm{w}_Z,\quad \forall w\in H^1_0(\dom).
	\end{equation}
	%
\end{lemma}	

\begin{lemma}
	\label{lem:aubinlionsmodified}
	Let  $\{\hu_h\}_{h>0}, \{u_h\}_{h>0}$ be  sequences of functions $\hu_h,u_h: [0,T]\times\dom\to \R^d$ satisfying
	\begin{equation}\label{eq:samelimit}
		\lim_{h\to 0}\norm{ u_h-\hu_h}_{L^2([0,T]\times \dom)}=0,	
		\end{equation}
	and
	\begin{equation}
		\label{eq:aprioriAL}\begin{split}
			& u_h\in L^2(0,T;L^2(\dom)\cap \Vh),\\
			&\hu_h\in  L^2(0,T;H^1_0(\dom)),\\
			&\partial_t \hu_h\in L^2(0,T;(H^k(\dom)\cap H^1_{\Div}(\dom))^*),	
		\end{split}
	\end{equation}
	uniformly in $h>0$ where $k\in \N$. Then the sequences $\{\hu_h\}_{h>0}$ and $\{u_h\}_{h>0}$ are precompact in $L^2([0,T]\times\dom)$.
\end{lemma}

\begin{proof}
	The proof follows the proof of the Aubin-Lions-Simon lemma with modifications to account for the fact that $\partial_t \hu_h\in (H^k(\dom)\cap H^1_{\Div}(\dom))^*$ but $\hu_h$ is not divergence free.
	
	From the Banach-Alaoglu theorem, we obtain that $u_h\weak u$ and $\hu_h\weak \hu$ in $L^2([0,T]\times\dom)$ up to a subsequence, for the ease of notation still denoted by $h>0$. From Lemma~\ref{lem:divfree2}, we obtain that the limit $u$ is divergence free (and $\PVh u =u$), and from Lemma~\ref{lem:samelimits}, we obtain that $u=\hu\in L^2(0,T;L^2_{\Div}(\dom))\cap L^2(0,T;H^1_0(\dom))$ almost everywhere and in $L^2([0,T]\times\dom)$. We let $v_h:=\hu_h-u$ and $w_h=u_h-u$, then $v_h$ and $w_h$ satisfy~\eqref{eq:samelimit} and~\eqref{eq:aprioriAL} and converge weakly to $0$ in $L^2([0,T]\times\dom)$. Note also that $w_h = \PVh w_h$.
	Let us denote $Z =  (H^k(\dom)\cap H^1_{\Div}(\dom))^*$. We start by showing that for a.e. $t\in [0,T]$,
		\begin{equation}\label{eq:uniformZ}
		\norm{v_h(t)}_{Z}\leq C,
	\end{equation}
	and that
	\begin{equation}\label{eq:aeliminZ}
		\norm{v_h(t)}_Z \stackrel{h\to 0}{\longrightarrow} 0,\quad \text{a.e. }\, t\in [0,T].
	\end{equation}
	Then, by the Lebesgue dominated convergence theorem, 
	we can conclude that $v_n\to 0$ in $L^2(0,T;Z)$. Without loss of generality, we let $t=0$. Then for any $s,t>0$,
	\begin{equation*}
		\begin{split}
			v_h(0)&= v_h(t)-\int_0^t v'_h(\tau)d\tau\\
			& = \frac{1}{s}\left(\int_0^s v_h(t)dt -\int_0^s \int_0^t v_h'(\tau)d\tau dt\right)\\
			&:= a_h + b_h.	
	\end{split}\end{equation*}
	We can rewrite $b_h$ for all $h>0$ as
	\begin{equation*}
		b_h = -\frac{1}{s}\int_0^s (s-t) v_h'(t)dt.
	\end{equation*}
	Thus given any $\epsilon>0$, we can choose $s>0$ such that for all $n\in\N$ it holds
	\begin{equation*}
		\norm{b_h}_{Z}\leq \int_0^s \norm{v_h'(t)}_{Z} dt\leq \sqrt{s}\norm{v_h'}_{L^2(0,T;Z)} \leq \frac{\epsilon}{2},
	\end{equation*}
	since $v_h'\in L^2(0,T;Z)$ uniformly in $h>0$ by the assumptions.
	On the other hand, we have, since $Z\subset (H^1_0(\dom)\cap H^k(\dom))^*$ that
	\begin{equation*}
		\norm{a_h}_{Z}\leq C \norm{a_h}_{(H^1_0(\dom)\cap H^k(\dom))^*} \leq  C \norm{a_h}_{L^2(\dom)}= C \norm{\frac{1}{s}\int_0^s v_h(t)dt}_{L^2(\dom)}.%\\
	\end{equation*}
	Due to the compact embedding $H^1_0(\dom)\subset\subset L^2(\dom)$ and~\eqref{eq:aprioriAL}, we have that $\frac{1}{s}\int_0^s v_h(t)dt\in H^1(\dom)$ is precompact in $L^2(\dom)$ and converges up to subsequence to zero in $L^2(\dom)$.
	Thus, we can find $\bar{h}>0$ such that for $h<\bar{h}$ (along a subsequence),
	\begin{equation*}
		\norm{a_h}_{Z}<\frac{\epsilon}{2}.
	\end{equation*}
	Thus~\eqref{eq:aeliminZ} holds.
	The uniform bound~\eqref{eq:uniformZ} follows in the same way, going over these estimates and observing that $a_h$ and $b_h$ are uniformly bounded independently of the time $t\in [0,T]$. 
	Hence, we obtain that $v_h(t)\to 0$ almost everywhere in $[0,T]$ and hence in $L^2(0,T;Z)$.
	Next, we have by Lemma~\ref{lem:lionslike},  that for any $\eta>0$, there exists $C_{\eta}>0$ and $h_\eta>0$ such that for all $h<h_\eta$,
	\begin{equation*}
		\norm{\PVh v_h}_{L^2([0,T]\times\dom)}\leq \eta \norm{v_h}_{L^2(0,T;H^1_0(\dom))}+C_{\eta}\norm{v_h}_{L^2(0,T;Z)}.
	\end{equation*}
	By the assumption~\eqref{eq:aprioriAL}, we have that
	\begin{equation*}
		\norm{v_h}_{L^2(0,T;H^1_0(\dom))}\leq C_0,
	\end{equation*}
	for some constant $C_0>0$ independent of $h$. 
	We now let $\epsilon>0$ arbitrary and pick $\eta = \epsilon/(2 C_0)$ (and $h_\eta$ correspondingly). As we have already shown that $\norm{v_h}_{L^2(0,T;Z)}\to 0$, we can choose $0<h_1\leq h_\eta$ such that for all $h<h_1$, we have
	\begin{equation*}
		C_\eta\norm{v_h}_{L^2(0,T;Z)}\leq \frac{\epsilon}{2}.
	\end{equation*}
	Thus for $h<h_1$, we have
	\begin{equation}\label{eq:PVhvconv}
		\norm{\PVh v_h}_{L^2([0,T]\times\dom)}\leq \epsilon.
	\end{equation}
	So we obtain that $\PVh v_h \to 0$ strongly in $L^2([0,T]\times\dom)$ along a subsequence. Next, we write
	\begin{equation*}
		v_h = v_h - \PVh v_h -w_h +\PVh w_h +\PVh v_h,
	\end{equation*}
	using that $w_h = \PVh w_h$. So we can estimate using the triangle inequality,
	\begin{align*}
		\norm{v_h}_{L^2([0,T]\times\dom)}&\leq \norm{v_h-w_h}_{L^2([0,T]\times\dom)} + \norm{\PVh v_h-\PVh w_h}_{L^2([0,T]\times\dom)} + \norm{\PVh v_h}_{L^2([0,T]\times\dom)}\\
		&\leq 2\norm{v_h-w_h}_{L^2([0,T]\times\dom)}   + \norm{\PVh v_h}_{L^2([0,T]\times\dom)},
	\end{align*}
	using that $\PVh$ is an orthogonal projection. The first term on the right hand side goes to zero by assumption~\eqref{eq:samelimit} and the second goes to zero thanks to~\eqref{eq:PVhvconv}. Thus $v_h\to 0$ in $L^2([0,T]\times\dom)$ up to a subsequence and by~\eqref{eq:samelimit} again, also $w_h\to 0$ in $L^2([0,T]\times\dom)$.
	Thus $v_h, \PVh v_h,w_h$ are precompact in $L^2([0,T]\times\dom)$ implying the precompactness of the sequences $\{\hu_h\}_{h>0}$, $\{u_h\}_{h>0}$ and $\{\PVh \hu_h\}_{h>0}$ in $L^2([0,T]\times\dom)$.
\end{proof}

\end{document}
